\documentclass[11pt]{article}
\usepackage[paper=a4paper,left=24mm,right=24mm,top=20mm,bottom=20mm]{geometry}
\usepackage[slantfont,boldfont]{xeCJK}
\usepackage{amsmath}                                % For textit, textbf, etc
\usepackage{graphicx}                               % For includegraphics
\usepackage{wrapfig}                                % For wrapfig environment
\usepackage{paralist}                               % For compactitem 
%\linespread{1.5}
\setmainfont{Times New Roman}	%缺省英文字体.serif是有衬线字体sans serif无衬线字体
\setCJKmainfont{KaiTi}	% 设置正文罗马族的CJK字体，影响\rmfamily和\textrm 的字体
\setCJKsansfont{SimSun}	% 设置正文无衬线族的CJK字体，影响\sffamily和\textsf 的字体
\setCJKmonofont{FangSong}	% 设置正文等宽族的CJK字体，影响\ttfamily 和 \texttt 的字体

%------------------------------设置字体大小------------------------%
\newcommand{\chuhao}{\fontsize{42pt}{\baselineskip}\selectfont}     %初号
\newcommand{\xiaochuhao}{\fontsize{36pt}{\baselineskip}\selectfont} %小初号
\newcommand{\yihao}{\fontsize{28pt}{\baselineskip}\selectfont}      %一号
\newcommand{\erhao}{\fontsize{21pt}{\baselineskip}\selectfont}      %二号
\newcommand{\xiaoerhao}{\fontsize{18pt}{\baselineskip}\selectfont}  %小二号
\newcommand{\sanhao}{\fontsize{15.75pt}{\baselineskip}\selectfont}  %三号
\newcommand{\sihao}{\fontsize{14pt}{\baselineskip}\selectfont}         %四号
\newcommand{\xiaosihao}{\fontsize{12pt}{\baselineskip}\selectfont}  %小四号
\newcommand{\wuhao}{\fontsize{10.5pt}{\baselineskip}\selectfont}    %五号
\newcommand{\subwuhao}{\fontsize{10pt}{\baselineskip}\selectfont}    %次五号
\newcommand{\xiaowuhao}{\fontsize{9pt}{\baselineskip}\selectfont}   %小五号
\newcommand{\liuhao}{\fontsize{7.875pt}{\baselineskip}\selectfont}  %六号
\newcommand{\qihao}{\fontsize{5.25pt}{\baselineskip}\selectfont}    %七号

\pagestyle{empty}                                   % No pagenumbers/headers/footers

%%% Custom sectioning (sectsty package)
\usepackage{sectsty}
\sectionfont{%			                            % Change font of \section command		                
	\sectionrule{0pt}{0pt}{-4pt}{1pt}
}
\subsectionfont{%			                        % Change font of \subsection command	                
	%\sectionrule{0pt}{0pt}{-10pt}{1pt}
}

%%% Macros
\newcommand{\newpart}[1]{
	\vspace*{-0.3cm}
	\section*{#1}
}
\newcommand{\newsubpart}[1]{\subsection*{#1}}

\newcommand{\name}[1]{
    \begin{center}
		\yihao #1       % Name
	\end{center}
}


\newcommand{\info}[1]{
	\noindent
	\vspace*{-1.0cm}
	\begin{center}
		#1
	\end{center}
}

\newcommand{\experience}[3]{
	\noindent
	\text{#1} \text{#2} \hfill \text{#3}
	\vspace{0.5em}
	\par
}

\newcommand{\education}[4]{
	\noindent
	\textbf{#1}                                     % School
	\hspace{0.1em}
	\text{#2}                                   	% Degree
	\hspace{0.1em}
	\text{#3} \hfill \text{#4}
	\par
}

%%% Begin Document
\begin{document}
	
%\begin{wrapfigure}{r}{0.3\textwidth}
%    \vspace*{-2.5cm}
%	\hfill\includegraphics[width=0.14\textwidth]{avatar.jpg}
%	\vspace{-1cm}
%\end{wrapfigure}

\MyName{林泽}


%%% Personal details

\BasicInfo{Tel: 19370575110}
\BasicInfo{Email: zee\_lin@foxmail.com}

%%% Education
\NewPart{教育经历}
\Postgraduate{2022.09 - 至今}
{华东师范大学}
{硕士在读}
{软件工程}
{预计于2025年06月毕业}

\Undergraduate{2018.09 - 2022.06}
{浙江工业大学}
{本科}
{计算机科学与技术}
{GPA： 3.54/5.0，综合排名： 5/112}
{曾获省政府奖学金、校优秀学生奖学金、校优秀本科毕业生等荣誉}

\vspace*{-0.2cm}

%%% Work experience
\NewPart{实习经历}
\Timestamp{蘑菇街 Weshop部门}{爬虫工程师}{2021.11 - 2022.02}
借助爬虫框架Scrapy，通过XPath定位的规则爬取国外电商平台中每个商品以Json格式存储的Sku数据。

%%% Skills and Qualifications
\NewPart{技能清单}
%\NewSubPart{JAVA, Python, LaTex}
熟练掌握JAVA, C++, Python, Latex等编程语言

%%% Awards
\NewPart{荣誉奖项}
\Timestamp{2020中国机器人大赛先进视觉赛--3D识别冠军（国家一等奖）}{}{2020.11}
\Timestamp{2020中国机器人大赛先进视觉赛--3D测量冠军（国家一等奖）}{}{2020.11}
\Timestamp{2021年第十六届全国大学生智能车竞赛--讯飞智能餐厅组别一等奖}{}{2021.08}
\Timestamp{第七届中国国际“互联网+”大学生创新创业大赛铜奖}{}{2021.10}

%%% Project experience
\NewPart{项目经历}
\Timestamp{基于计算机视觉的安全帽佩戴检测与预警系统}{}{2019.12 - 2020.05}
在Yolov3算法框架基础上进行剪枝、蒸馏等模型压缩操作并部署在本地服务器上，并向云端服务器实时更新报警信息。


\NewPart{其他}
\begin{itemize}	
	\item 英语水平： CET-4： 536，CET-6：516
	\item 个人博客： https://CompetitiveLin.github.io
\end{itemize}

\end{document}
